\chapter*{Rezumat}
\addcontentsline{toc}{chapter}{Rezumat}

Lucrarea de licen\c t\u a \emph{Aplica\c tie web pentru organizarea \c si gestionarea evenimentelor} propune dezvoltarea unei platforme complete destinate organizatorilor \c si participan\c tilor la conferin\c te, meeting-uri, concerte, evenimente culturale \c si, \^intr-o m\u asur\u a mai mic\u a, competi\c tii sportive. Aplica\c tia (implementare: proiectul ManFast) permite publicarea evenimentelor, filtrarea dup\u a nume, jude\c t \c si dat\u a, achizi\c tionarea biletelor nominale cu plat\u a online prin Stripe \c si administrarea participan\c tilor printr-un panou dedicat.

Solu\c tia tehnic\u a adopt\u a arhitectura client-server: frontend React 19 cu Material UI 9, backend Node.js/Express, baz\u a de date PostgreSQL. Autentificarea utilizeaz\u a token-uri JWT, parole hash-uite cu bcrypt, verificare email la \^inregistrare (cod de 6 cifre) \c si resetare parol\u a. Administratorii gestioneaz\u a evenimente (inclusiv imagini), efectueaz\u a check-in \c si promoveaz\u a al\c ti administratori.

Documentul prezint\u a motiva\c tia temei, starea domeniului, proiectarea \c si implementarea solu\c tiei, capturi din aplica\c tia func\c tional\u a \c si concluzii. Platforma reprezint\u a o solu\c tie scalabil\u a, adaptabil\u a contextului rom\^anesc (pl\u a\c ti RON, jude\c te).

\textbf{Cuvinte cheie}: aplicație web, gestionarea evenimentelor, bilete digitale, plată online, autentificare securizată.

\vspace{1cm}
\chapter*{Abstract}
\addcontentsline{toc}{chapter}{Abstract}

The bachelor's thesis Web application for organizing and managing events proposes the development of a complete platform for organizers and participants in conferences, meetings, concerts, cultural events, and to a lesser extent sports competitions. The application (implementation: ManFast project) enables event publishing, filtering by name, county and date, purchasing nominal tickets with online payment via Stripe, and participant management through a dedicated admin panel.

The technical solution adopts a client-server architecture: React 19 frontend with Material UI 9, Node.js/Express backend, PostgreSQL database. Authentication uses JWT tokens, bcrypt password hashing, email verification on registration (6-digit code), and password reset. Administrators manage events (including images), perform check-in, and promote other administrators.

The document presents the motivation, domain state, solution design and implementation, application screenshots, and conclusions. The platform represents a scalable solution adapted to the Romanian context (RON payments, counties).

\textbf{Keywords}: web application, event management, digital tickets, online payment, secure authentication.

\clearpage